\scanpage{118}\mspage{63}%
nation peut-être la plus jolie est par
\[
 \text{\uline{Parf}}\,S\longrightarrow\mathrm{Cat}_{/S}
 \tag*{$(*)$}
\]
qui donne le foncteur $i$ quand on travaille avec $W_S^d$, le foncteur $j$ quand on travaille avec $W_S^g$. On voit \inserted{\unsure{alors|donc}} tout de suite que si on compose avec le foncteur pertinent, soit $\unsure{\rho|p}$ soit $\sigma$, défini par l'inclusion $W_S^{\unsure{d|g}}\unsure{\in|\subset}W_S$ resp $W_S^g\subset W_S$ (comme parties de $\mathrm{Fl}(\mathrm{Cat}_{/S})$), on trouve le foncteur \deleted{canonique}
\[
 (\text{\uline{Parf}}\,S)W_S^{-1}\overset{\approx}{\unsure{\longrightarrow|\text{---}}}(\mathrm{Cat}_{/S})W_S^{-1}
 \tag*{$\unsure{(**)|(*)'|(*)}$}
\]
qui est l'\uline{équivalence canonique} entre ces deux catégories. Ainsi le \deleted{\unsure{des}} foncteur correspondant à $i,j$ page 51, exprimé en termes de $\mathrm{Cat}_{/S}$ comme foncteur en sens inverse de $\rho,\sigma$ (p.~62), est-il le composé de l'équiv. quasi-inverse de \unsure{$(**)$}, laquelle s'explicite par le foncteur \unsure{$\mathcal C$|$C$}, et du foncteur sur les \unsure{localisés|localis.} $W_S^d$ ou $W_S^g$ déduit de $(*)$, i.e.
\diagram{\scriptsize\begin{tikzcd}[ampersand replacement=\&,column sep=small,row sep=large]
 {} \& {} \& \mathrm{HOT}_W(S)\approx(\mathrm{Cat}_{/S})(W_S^d)^{-1}\\
 \mathrm{HOT}_W^{\unsure{\mathrm{lc}|\ell\mathrm{c}}}(S\unsure{.|{}})\approx(\mathrm{Cat}_{/S})W_S^{-1}
 \arrow[urr,"i"]\arrow[rrd,"j"']\arrow[r,"{\deleted{\unsure{\approx}}\mathcal C}"]
 \&(\text{\uline{Parf}}\,S)W_S^{-1}
 \arrow[ur,"{\text{car }W_S=W_S^d\text{ sur \uline{Parf}}}"']
 \arrow[dr,"{\text{car }W_S=W_S^g\text{ sur \uline{Parf}}}"]
 \&{}\\
 {} \& {} \& \mathrm{HOT}_W^\circ(S)\approx(\mathrm{Cat}_{/S})(W_S^g)^{-1}\,.
\end{tikzcd}}
% archive sheet number: 117
